\documentclass[nonewpage]{xourse}


% Start \part op een nieuwe pagina (Ximera doet dat niet per default)
\xpretocmd{\part}{\clearpage}{}{}

\newcommand{\richting}{Wetenschappen 5de jaar}%De klas/richting --> hoofding header
% \documentclass{xourse}
%in
%\title{\includegraphics[width=0.7\textwidth]{pic5/elektromagnetisme.jpg}\\[2cm]Fysica 5de jaar - Wetenschappen}
\title{Fysica 5de jaar - Wetenschappen}
\author{Daan Lipkens}
\begin{document}
\begin{abstract}
    Cursus Fysica 5de jaar Wetenschappen domeinoverschrijvende doorstroom.
\end{abstract}
\maketitle
%HOE MAAK IK DIT IN TITLE BLAD 
%\begin{figure}[h]
% \begin{center}
% \includegraphics[width=0.85\textwidth]{pic5/elektromagnetisme.jpg}
% \caption{De Joint European Torus (JET) is het grootste kernfusie-experiment ter wereld. Met behulp van krachtige elektromagnetische velden houdt deze tokamak het extreem hete plasma op zijn plaats, zodat de wandmaterialen niet smelten.}
% \end{center}
% \end{figure}
%\xmtitle{Fysica 5de jaar - Wetenschappen}{}
% \part voor delen
%\chapterstyle voor hoofdstukken >> \activity{hoofdstuk ....}
%\sectionstyle voor delen in het hoofdstuk >> \activity{met onderdelen}
% Inleiding, informatie, ...


% INFO
% toepassingen onder \begin{example}
% Quizjes als \question met \begin{selectAll}
% voorbeeld oefeningen met align x &= berekening && (uitleg)
% komma's met {,} in math mode


\part{Inleiding}
\chapterstyle
\activity{inleiding_stand_52.tex}
\sectionstyle
\activity{inleiding_praktische_richtlijnen_52.tex}
\activity{inleiding_hoe_fysica_te_studeren_5.tex}
\activity{inleiding_vectorkleuren.tex}


%--------ELEKTRICITEIT----------------
\part{Elektrische krachtwerking} %parts nog aanpassen
%COULOMBKRACHT of ELEKTRISCHE KRACHTWERKING
%picture path: pic5/elektriciteit/coulomb/

\chapterstyle

\activity{elektrische_krachtwerking_inleiding.tex} %inleiding hoofdstuk coulombkracht

\sectionstyle

\activity{elektrische_krachtwerking_lading.tex} %ladingen en soorten ladingen
\activity{elektrische_krachtwerking_tribo-elekrisch_effect.tex} %tribo-elektrisch effect
\activity{elektrische_krachtwerking_coulombkracht.tex} %coulombkracht met cosinusregel
\activity{elektrische_krachtwerking_materialen_en_geleiding.tex}
\activity{elektrische_krachtwerking_aantrekken_ongeladen_deeltjes.tex}
\activity{elektrische_krachtwerking_elektroscoop.tex}
\activity{elektrische_krachtwerking_opladen_inductie.tex}
\activity{elektrische_krachtwerking_samenvatting.tex}
%oefeningen
\chapterstyle
\activity{elektrische_krachtwerkingoef_inleiding.tex}
\sectionstyle
\activity{elektrische_krachtwerkingoef_vragen.tex}
\activity{elektrische_krachtwerkingoef_basis.tex}
\activity{elektrische_krachtwerkingoef_1lijn.tex}
\activity{elektrische_krachtwerkingoef_driehoek.tex}


\part{Elektrisch veld}
%ELEKTRISCH VELD
%picture path: pic5/elektriciteit/veld/
 
\chapterstyle
\activity{elektrischveld_inleiding.tex}
 
 
\sectionstyle
\activity{elektrischveld_concept.tex} %wat is E-veld, analogie zwaarte kracht en hoe kracht staat op een lading in een veld
\activity{elektrischveld_veldlijnen.tex} % veldlijnen en soorten velden
\activity{elektrischveld_puntlading.tex} % E-veld van puntlading
\activity{elektrischveld_dipool.tex} %E-veld tussen 2 ladingen
\activity{elektrischveld_homogeen_veld.tex} % homogeen veld
\activity{elektrischveld_materie_in_veld.tex} % geleider en isolator in veld, elektrische schermwerking + toepassingen
\activity{elektrischveld_toepassingen.tex}
\activity{elektrischveld_samenvatting.tex} %todo
 
% %oefeningen
\chapterstyle
\activity{elektrischveldoef_inleiding.tex}
\sectionstyle
\activity{elektrischveldoef_vragen.tex}
\activity{elektrischveldoef_basis.tex}
\activity{elektrischveldoef_1lijn.tex}
\activity{elektrischveldoef_driehoek.tex}



\part{Potentiële elektrische energie bij homogeen veld en spanning}
%POTENTIËLE ELEKTRISCHE ENERGIE BIJ HOMOGEEN VELD EN SPANNING
%picture path: pic5/elektriciteit/potentiaal/
\chapterstyle
\activity{potentiaal_inleiding.tex}

\sectionstyle
\activity{potentiaal_potentiele_energie.tex}
\activity{potentiaal_potentiaal.tex}
\activity{potentiaal_equipotentiaaloppervlakken.tex}
\activity{potentiaal_spanning.tex}
\activity{potentiaal_energie_arbeid.tex}
\activity{potentiaal_samenvatting.tex}

\chapterstyle
\activity{potentiaaloef_inleiding.tex}
\sectionstyle
\activity{potentiaaloef_vragen.tex}
\activity{potentiaaloef_basis.tex}
\activity{potentiaaloef_energie.tex}

% %--------MAGNETISME----------------
 \part{Elektromagnetisme}
%HET MAGNETISCH VELD (permanente magneten, rechte geleider, spoel, elektromagneet)
%figuren: TikZ in de bestanden zelf
\chapterstyle
\activity{magnetisme_inleiding.tex}

\sectionstyle
\activity{magnetisme_permanente_magneten.tex} %polen, influentie, ferromagnetisme, Extra: Weissgebieden
\activity{magnetisme_veld.tex} %B-vector, veldlijnen, homogeen veld, aardmagnetisme, analogie velden
\activity{magnetisme_rechte_geleider.tex} %Oersted, rechterhandregel, B = mu0 I/(2 pi r), superpositie
\activity{magnetisme_spoel.tex} %winding, spoel, B = mu0 N I/l, bewegende ladingen als bron
\activity{magnetisme_elektromagneet.tex} %ijzeren kern, mu_r, elektromagneet, Extra: toepassingen
\activity{magnetisme_samenvatting.tex}

\chapterstyle
\activity{magnetismeoef_inleiding.tex}
\sectionstyle
\activity{magnetismeoef_vragen.tex}
\activity{magnetismeoef_basis.tex}
\activity{magnetismeoef_gevorderd.tex}

 \part{Lorentzkracht}
 \part{Inductiespanning}

% \newpage
% %--------KERNFYSICA----------------
 \part{Deeltjesfysica}\label{hoofdstuk:deeltjesfysica}
 \part{De atoomkern} %deeltjesfysica apart deel?
 \part{Radioactiviteit}
 \part{Dosis}


%-----------AI----------------
\part{AI-assistent voor extra oefeningen}
\chapterstyle
\activity{AI_inleiding.tex}

%\sectionstyle
\activity{AI_elektrische_krachtwerking_2.tex}
\activity{AI_elektrischveld_2.tex}
\activity{AI_potentiaal_2.tex}

%---------APPENDIX---------------

\part{Appendix} %nummering anders

\chapterstyle
\activity{appendix_vectoren_inleiding.tex}

\sectionstyle
\activity{appendix_vectoren_scalar.tex}
\activity{appendix_vectoren_component.tex}
\activity{appendix_vectoren_som.tex}
\activity{appendix_vectoren_verschil.tex}
\activity{appendix_vectoren_scalaire_vermenigvuldiging.tex}
\activity{appendix_vectoren_vectorieel_product.tex}


\chapterstyle
\activity{appendix_standaardvorm_inleiding.tex}

\sectionstyle
\activity{appendix_standaardvorm_si.tex}
\activity{appendix_standaardvorm_omzetten_eenheden.tex}
\activity{appendix_standaardvorm_bc.tex}
\activity{appendix_standaardvorm_oefeningen.tex}


\chapterstyle
\activity{appendix_energie_inleiding.tex}

\sectionstyle
\activity{appendix_energie_arbeid.tex}
\activity{appendix_energie_vormen.tex}
\activity{appendix_energie_behoud.tex}


\chapterstyle
\activity{appendix_gegevens_constanten.tex}

%Vector rekenen
    %vectoren
    %cosinus regel

%nuclide tabel
%standaard vorm
%logaritme en exponentiële functies
%---------Bibliografie
\part{Bibliografie}
\activity{bibliografie_52.tex}
\end{document}
